\scanpage{93}\mspage{46}
$\overline{\underline{\mathrm{VI}}}$ Je commencerai par la question e)
de la page précédente, qui avait déjà été
laissée en suspens au \S\,$\unsure{\mathrm{IV}|\mathrm{III}|\mathrm{VI}}$.

\deleted{\uline{Proposition}} Soit $\Lambda$ un anneau quelconque, \inserted{$M$ $\Lambda$-mod libre de rg 2}
considérons l'action de $\mathrm{Sl}_\Lambda(M)$ sur $M^*=\widetilde S_M$
de $\mathrm{Sl}_\Lambda(M)'$ sur $M^*/\pm1=S_M$. On cherche
les commutants de ces actions -- \deleted{dans les groupes $\unsure{\circlearrowright|\mathfrak S}_M$} \deleted{avec}
\inserted{dans les groupes des permutations de $\widetilde S_M$ resp. $\circlearrowright_M$}
présomption que c'est l'image de $\Lambda^*$
resp. de $Z'_\Lambda=\Lambda^*/\pm1$. Comme $M^*$ et a
fortiori \inserted{son quotient} $S_M$ \deleted{est} sont des espaces homogènes
sous \inserted{le Groupe d'\unsure{opération|opérateurs} $G$} $\mathrm{Sl}_\Lambda(M)$ resp. $\mathrm{Sl}_\Lambda(M)'$, le commutant
\inserted{dans $\circlearrowright_S$ de $G$} est isomorphe à $N_G(H)/H$, où $H$ est
le stabilisateur d'un point $s$ de $S$, de
façon précise on aura un \deleted{homomorphisme}
\deleted{\unsure{su|iso}}\inserted{isom}
\[
\deleted{\unsure{C}}Z\xrightarrow{\sim}N_G(H)/H
\]
par la relation entre $\deleted{\unsure{C}}Z$ et $N_G(H)$
définie par la formule
\[
u(s)=g\cdot s\qquad u\in Z,\ g\in N_G(H)\tag*{(48)}
\]

Calculons donc le normalisateur de
$H$, dans les deux cas, en prenant une base
$(x,y)$ de $M$, et en prenant $s=x$ ou $s=\dot x$.
% archive sheet number: 92
